\documentclass[11pt, a4paper]{article}

% ---------- PACKAGES ----------
\usepackage[top=1.2cm, bottom=1.2cm, left=1.5cm, right=1.5cm]{geometry}
\usepackage{titlesec}
\usepackage{enumitem}
\usepackage{hyperref}
\usepackage{xcolor}
\usepackage{fontspec}
\usepackage{microtype}

\usepackage[normalem]{ulem}  % For underline
\setmainfont{Times New Roman}

% ---------- COLORS (for links only) ----------
\hypersetup{
  colorlinks=true,
  urlcolor=blue
}

\pagestyle{empty}
\setlength{\parindent}{0pt}
\setlength{\parskip}{1.5pt}
\setlength{\emergencystretch}{3em}  % Avoid overfull hbox
\definecolor{sectiongray}{RGB}{60,60,60}


% ---------- SECTION STYLE ----------
\definecolor{sectiongray}{RGB}{60,60,60}

\titleformat{\section}
  {\large\bfseries\color{sectiongray}}
  {}
  {0em}
  {}
  [\vspace{2pt}\titlerule]

\titlespacing{\section}{0pt}{10pt}{6pt}



% ---------- LIST STYLE ----------
\setlist[itemize]{
    leftmargin=1.2em,
    itemsep=1.5pt,
    topsep=2pt,
    parsep=0pt
}

% ---------- DOCUMENT ----------
\begin{document}

% ================= HEADER =================
\begin{center}
{\LARGE \textbf{BRYAN STEVEN MENOSCAL SANTANA}}
\end{center}

\vspace{6pt}

\begin{tabular*}{\textwidth}{@{\extracolsep{\fill}} l l}
\textbf{Ubicación:} Ecuador, Manabí, Manta & \textbf{Teléfono:} +593 99 739 9441 \\
\textbf{Correo:} \href{mailto:bryanmenoscal2005@gmail.com}{bryanmenoscal2005@gmail.com}
& \textbf{Portafolio:} \href{https://stevsant.vercel.app?ref=cv-es}{stevsant.vercel.app} \\
\textbf{LinkedIn:} \href{https://linkedin.com/in/bryanmenoscal26}{linkedin.com/in/bryanmenoscal26}
& \textbf{GitHub:} \href{https://github.com/StevSant}{github.com/StevSant} \\
\end{tabular*}

\vspace{6pt}
\hrule
\vspace{8pt}


% ================= RESUMEN =================
\section*{RESUMEN}

Desarrollador de Software en formación con enfoque Fullstack (Python/Angular).
Ganador del 3er lugar en la Hackathon Viamatica 2025 (entre 450+ participantes).
Apasionado por la arquitectura de software escalable y la Inteligencia Artificial.
Inglés avanzado (C1) y autodidacta, con capacidad demostrada para liderar iniciativas técnicas y resolver problemas bajo presión.

% ================= IDIOMAS =================
\section*{IDIOMAS}
\vspace{-2pt}
\begin{tabular}{p{0.25\textwidth} p{0.7\textwidth}}
\textbf{Español:} & Nativo \\
\textbf{Inglés:} & C1 (Avanzado)
\end{tabular}
\vspace{6pt}
\hrule
\vspace{8pt}

% ================= SKILLS =================
\section*{SKILLS}
\vspace{-2pt}
\begin{tabular}{p{0.25\textwidth} p{0.7\textwidth}}
\textbf{Backend:} & Django, Django REST Framework, Flask, FastAPI, NestJS \\
\textbf{Frontend:} & Angular (principal), React (básico) \\
\textbf{Móviles:} & Kotlin (básico) \\
\textbf{Lenguajes:} & Python, JavaScript, TypeScript, Java/Kotlin (básico) \\
\textbf{Bases de datos:} & MySQL, PostgreSQL, MongoDB \\
\textbf{Herramientas:} & Git, Docker, GitHub/GitLab, Postman, VS Code \\
\textbf{Cloud:} & Vercel, Azure, Google Cloud Run \\
\textbf{Metodologías:} & Scrum, desarrollo ágil, integración continua \\
\end{tabular}
\vspace{6pt}
\hrule
\vspace{8pt}

% ================= PROYECTOS =================
\section*{PROYECTOS Y EXPERIENCIA RELEVANTE}

\textbf{AGENTE DE EVALUACIÓN DE RIESGO PARA PYMES --- HACKATÓN VIAMATICA (2025)} \hfill Ecuador, Guayas, Guayaquil \\
\textit{Backend \& AI Developer | Hackathon Viamatica 2025 (460+ participantes)} \hfill 2025

\begin{itemize}
\item \textbf{Logro:} 3er lugar entre 460+ participantes (equipos de 3), desarrollando un MVP funcional en menos de 3 días.
\item \textbf{Desarrollo:} Construí el agente de IA (LangChain) y la API REST (FastAPI) para evaluación de riesgo crediticio, implementando web scraping automatizado con Playwright para extraer datos de la Superintendencia de Compañías, el SRI y redes sociales (Twitter, Facebook, Instagram, LinkedIn).
\item \textbf{Impacto:} Introduje un score de reputación digital para PYMES, permitiendo evaluaciones crediticias más justas al incorporar datos de presencia online que normalmente se ignoran en métodos tradicionales.
\end{itemize}

\vspace{4pt}

\textbf{STREAMFLOWMUSIC} \hfill Proyecto Personal \\
\textit{Desarrollador Backend} \hfill 2025

\begin{itemize}
\item Diseñé y desarrollé una API RESTful escalable tipo Spotify utilizando Django REST Framework y Supabase.
\item Implementé un sistema de ingesta de contenido integrando la API de YouTube para automatizar la descarga y almacenamiento en la nube.
\item \textbf{Pruebas de carga:} Validé la robustez del sistema simulando 100 usuarios concurrentes, optimizando las consultas a base de datos para mantener baja latencia.
\end{itemize}

% ================= EDUCACION =================
\section*{EDUCACIÓN}

\textbf{UNIVERSIDAD LAICA ELOY ALFARO DE MANABÍ} \hfill Ecuador, Manabí, Manta \\
\textit{Ingeniería en Software} \hfill 2023 -- Presente

\begin{itemize}
\item Desarrollo de proyectos académicos y personales utilizando Django REST, Flask, FastAPI, NestJS y Angular, enfocados en resolver problemas reales con soluciones escalables.
\item Formación complementaria autodidacta en backend, frontend y metodologías ágiles (Scrum/Kanban), reforzando la práctica con experiencias aplicadas.
\end{itemize}

% ================= LIDERAZGO =================
\section*{LIDERAZGO Y VOLUNTARIADO}

\textbf{Club de Inteligencia Artificial ULEAM} \hfill Ecuador, Manabí, Manta \\
\textit{Encargado de gestión de contenidos} \hfill 2024 -- Presente

\begin{itemize}
\item Planificación de temarios y contenidos de formación en Inteligencia Artificial para la comunidad universitaria.
\item Apoyo en la organización de charlas y talleres tecnológicos, fomentando el aprendizaje colaborativo.
\end{itemize}

\end{document}
